% Self-contained proof module; needs no modular representation theory.
\section{An obstruction in every positive characteristic}

An exterior-power identity makes characteristic zero a necessary
condition, including for coefficient rings of characteristic $p^a$.

\begin{proposition}\label{prop:positive-characteristic}
Let $A$ be a nonzero commutative ring, and suppose that $m\cdot1_A=0$
for some integer $m\geq2$.  Let $P_m$ be the $m$th power-sum
symmetric function, regarded as an integral polynomial in the
complete symmetric functions $h_1,\ldots,h_m$.  Then
\[
 P_m(\rho_1,\ldots,\rho_m)=0
 \quad\text{in }K_0(\operatorname{Lat}_A(S_m)).
\]
In particular, the trivial classes cannot freely generate over any
nonzero local ring of positive characteristic.
\end{proposition}

\begin{proof}
For each $r\geq0$, write $\epsilon_r$ for the sign lattice of $S_r$,
with $\epsilon_0=\rho_0=1$.  For the permutation lattice
$A^r$, contraction by the sum of coordinates makes its exterior
algebra an exact complex
\[
 0\longrightarrow\bigwedge^r A^r\longrightarrow\cdots
 \longrightarrow\bigwedge^1 A^r\longrightarrow A
 \longrightarrow0.
\]
This complex is $S_r$-equivariant and split exact over $A$: contraction
has a contracting homotopy given by exterior multiplication by a
basis vector.  Its cycles are finite free, as can also be seen by
splitting $A^r$ into the kernel of the sum map and that basis line.
Since
\[
 \bigwedge^j A^r\simeq
 \operatorname{Ind}_{S_j\times S_{r-j}}^{S_r}
 (\epsilon_j\boxtimes\rho_{r-j}),
\]
its Euler relation is
\[
 \sum_{j=0}^r(-1)^j\rho_{r-j}\epsilon_j=0
 \qquad(r>0).
\]
Thus, if $H(t)=\sum_{r\geq0}\rho_rt^r$ and
$E(t)=\sum_{r\geq0}\epsilon_rt^r$, then $H(t)E(-t)=1$.
In particular $\epsilon_j$ is the usual elementary symmetric
polynomial in the $\rho_i$.

Now set $U=A^m$ with its permutation action, and let
$V=\ker(U\xrightarrow{\sum} A)$.  This is free of rank $m-1$,
with basis $e_i-e_m$ for $1\leq i<m$.  The invariant vector
$v=e_1+\cdots+e_m$ belongs to $V$ because $m=0$ in $A$.
Moreover,
\[
 v=\sum_{i=1}^{m-1}(e_i-e_m),
\]
so it is primitive and $W=V/Av$ is finite free of rank $m-2$.
The exterior-power filtration of
$0\to Av\to V\to W\to0$ gives
\[
 [\textstyle\bigwedge^i V]
 =[\textstyle\bigwedge^i W]+[\textstyle\bigwedge^{i-1}W].
\]
With exterior powers outside the possible range set to zero,
alternation gives
\begin{equation}\label{eq:augmentation-exterior-zero}
 \sum_{i=0}^{m-1}(-1)^i[\textstyle\bigwedge^iV]=0.
\end{equation}
On the other hand, $0\to V\to U\to A\to0$ yields
\[
 [\textstyle\bigwedge^iV]
 =\sum_{j=0}^i(-1)^{i-j}[\textstyle\bigwedge^jU].
\]
Substitution into~\eqref{eq:augmentation-exterior-zero} gives
\[
 0=\sum_{j=0}^{m-1}(m-j)(-1)^j
       \rho_{m-j}\epsilon_j.
\]
The right side is the coefficient of $t^m$ in
$tH'(t)E(-t)=tH'(t)/H(t)$, hence is
$P_m(\rho_1,\ldots,\rho_m)$ by the Newton identity.  The
polynomial $P_m$ is nonzero over $\mathbb Z$: its coefficient of
$h_m$ is $m$.  This proves the obstruction.
\end{proof}

\begin{example}
In characteristic $2$, the identity is
$2\rho_2-\rho_1^2=0$.
In characteristic $4$, the argument gives the nonzero polynomial
relation
\[
 4\rho_4-4\rho_1\rho_3-2\rho_2^2
       +4\rho_1^2\rho_2-\rho_1^4=0.
\]
The second example does not require a computation of the lattice
Grothendieck group of $S_2$ over $\mathbb Z/4\mathbb Z$.
\end{example}
